% Bagian: computability
% Bab: computability-theory
% Subbagian: complement-ce

\documentclass[../../../include/open-logic-section]{subfiles}

\begin{document}

\olfileid[id]{cmp}{thy}{cmp}
\olsection{Himpunan Terenumerasi Secara Komputabel Tidak Tertutup terhadap Komplemen}

Andaikan $A$ terenumerasi secara komputabel. Apakah komplemen~$A$,
$\Complement{A} = \Nat \setminus A$, selalu terenumerasi secara komputabel
juga? Teorema dan akibat berikut menunjukkan bahwa jawabannya adalah ``tidak.''

\begin{thm}
\ollabel{thm:ce-comp}
Misalkan $A$ sembarang himpunan bilangan asli. Maka $A$ dapat dihitung jika dan
hanya jika $A$ maupun $\Complement{A}$ terenumerasi secara komputabel.
\end{thm}

\begin{proof}
Arah maju mudah: jika $A$ dapat dihitung, maka $\Complement{A}$ juga dapat
dihitung ($\Char{A} = 1 \tsub \Char{\Complement{A}}$), sehingga keduanya
terenumerasi secara komputabel.

Untuk arah sebaliknya, andaikan $A$ dan~$\Complement{A}$ keduanya terenumerasi
secara komputabel. Misalkan $A$ merupakan domain~$\cfind{d}$ dan
$\Complement{A}$ merupakan domain~$\cfind{e}$. Definisikan $h$ dengan
\[
h(x) = \umin{s}{(T(d,x,s) \lor T(e,x,s))}.
\]
Dengan kata lain, pada masukan~$x$, $h$ mencari komputasi berhenti
$\cfind{d}$ atau komputasi berhenti $\cfind{e}$. Jika $x \in A$, pencarian
berhasil dalam kasus pertama; jika $x \in \Complement{A}$, pencarian berhasil
dalam kasus kedua. Jadi, $h$ merupakan fungsi dapat dihitung total. Sekarang,
untuk setiap~$x$, kita mempunyai $x \in A$ jika dan hanya jika
$T(d, x, h(x))$, yakni jika $\cfind{d}$ yang terdefinisi. Karena
$T(d, x, h(x))$ merupakan relasi dapat dihitung, $A$ dapat dihitung.
\end{proof}

\begin{explain}
Lebih mudah memahami yang terjadi dalam istilah komputasional informal: untuk
memutuskan $A$ pada masukan $x$, carilah komputasi berhenti $\cfind{d}$ dan
$\cfind{e}$. Salah satunya pasti berhenti; jika yang berhenti adalah
$\cfind{d}$, maka $x$ berada dalam~$A$, dan jika tidak, $x$ berada dalam
$\Complement{A}$.
\end{explain}

\begin{cor}
\ollabel{cor:comp-k}
$\Complement{K_0}$ tidak terenumerasi secara komputabel.
\end{cor}

\begin{proof}
Kita mengetahui bahwa $K_0$ terenumerasi secara komputabel, tetapi tidak dapat
dihitung. Jika $\Complement{K_0}$ terenumerasi secara komputabel, maka $K_0$
dapat dihitung menurut \olref{thm:ce-comp}, yang bertentangan dengan
\olref[ncp]{thm:K-0}.
\end{proof}

\end{document}

